\documentclass[proposal,indonesia]{ugmskripsi}
% Options:
%   skripsi / proposal        (default: skripsi)
%   satuspasi / duaspasi / satusetengahspasi (default: satusetengahspasi)
%   indonesia / inggris       (default: indonesia)

%------------------------------------------------------------
% Metadata Proposal
%------------------------------------------------------------
% TODO: Ganti metadata di bawah ini dengan data Anda sendiri
\titleind{Perancangan dan Implementasi Arsitektur Software as a Service untuk Sistem Ujian Daring dengan Integrasi Anti-Kecurangan Multi-Platform dan Penilaian Berbasis Item Response Theory}
% \titleeng{Design and Implementation of a Software as a Service Architecture for an Online Examination System with Multi-Platform Anti-Cheat Integration and Item Response Theory-Based Assessment}
\fullname{Nugroho Adi Susanto}
\idnum{23/520312/PA/22367}
\degree{BACHELOR}
\gelar{SARJANA}
\yearsubmit{2026}
\program{ILMU KOMPUTER}
\dept{ILMU KOMPUTER DAN ELEKTRONIKA INSTRUMENTASI}
\firstsupervisor{ILMU KOMPUTER}
\firstexaminer{NAMA DOSEN PENGUJI 1}
\secondexaminer{NAMA DOSEN PENGUJI 2}

%------------------------------------------------------------
% Packages
%------------------------------------------------------------
\usepackage[utf8]{inputenc}
\usepackage{graphicx}
\graphicspath{{figures/}}
\usepackage{amsmath}
\usepackage{amssymb}
\usepackage{booktabs}
\usepackage{longtable}
\usepackage{tabularx}
\usepackage{multirow}
\usepackage{float}
\usepackage{listings}
\usepackage{xcolor}
\usepackage{hyperref}
\usepackage{url}
\usepackage{enumitem}

% Code listing style
\lstset{
    basicstyle=\ttfamily\small,
    breaklines=true,
    frame=single,
    numbers=left,
    numberstyle=\tiny,
    keywordstyle=\color{blue},
    commentstyle=\color{gray},
    stringstyle=\color{orange},
    showstringspaces=false,
    tabsize=2
}

% Define YAML language for listings
\lstdefinelanguage{yaml}{
  keywords={true,false,null,y,n},
  sensitive=false,
  comment=[l]{\#},
  morecomment=[s]{/*}{*/},
  morestring=[b]',
  morestring=[b]"
}

% Define Docker language for listings
\lstdefinelanguage{docker}{
  alsoletter={-},
  keywords={FROM,RUN,CMD,LABEL,MAINTAINER,EXPOSE,ENV,ADD,COPY,ENTRYPOINT,VOLUME,USER,WORKDIR,ARG,ONBUILD,STOPSIGNAL,HEALTHCHECK,SHELL},
  sensitive=false,
  comment=[l]{\#},
  morestring=[b]',
  morestring=[b]"
}

% Set Logo size
\renewcommand{\logowidth}{5.5cm} 

%------------------------------------------------------------
% Document Body
%------------------------------------------------------------
\begin{document}

% Preliminary Pages (Proposal: minimal)
\cover
\titlepageind
\approvalpage

\tableofcontents

\abstractind
% Intisari dalam Bahasa Indonesia (150-250 kata)

Indonesia memiliki lebih dari 316 ribu sekolah yang terkoneksi internet\cite{indonesiabaikRibuSekolah}. Dari seluruh sekolah tersebut,
seluruhnya diwajibkan untuk mengikuti Tes Kemampuan Akademik (TKA) yang berskala nasional dan dilakukan
menggunakan sarana komputer atau \emph{Computer-Based Testing} (CBT). Hal ini menimbulkan kebutuhan akan 
kegiatan persiapan ujian TKA bagi tiap sekolah. Kegiatan persiapan ujian TKA harus mensimulasikan kondisi 
ujian TKA yang sesungguhnya, hal ini termasuk akan penggunaan software ujian daring yang mirip.

Platform persiapan ujian daring (Computer-Based Testing/CBT) saat ini menghadapi sejumlah tantangan utama untuk digunakan 
sebagai software persiapan atau simulasi bagi institusi. tiga tantangan utama tersebut yaitu integritas ujian dari kecurangan 
peserta, ketersediaan algoritma penilaian IRT pada software dan kapasitas software untuk menampung dan mengelola data 
\emph{concurrent} submision dalam jumlah besar atau lebih dari 1000 request per detik. Jumlah data sangat penting bagi 
institusi karena keakuratan parameter IRT sangat bergantung pada jumlah data yang tersedia. 
Penyimpangan statistik dari jumlah data yang tersedia akan mempengaruhi keakuratan parameter IRT sehingga menyebabkan
penurunan akurasi skor atau deviasi skor tryout terhadap ujian sebenarnya dan mempersulit institusi untuk 
menilai dan mempersiapkan siswa.

Penelitian ini bertujuan merancang dan mengimplementasikan arsitektur SAAS multi-tenant bagi lebih dari satu institusi 
sebagai platform simulasi ujian daring berbasiskan IRT dengan kemampuan deteksi kecurangan untuk lebih dari 1000 
peserta ujian konkurren. Sistem anti-kecurangan akan menerapkan deteksi berbasis perilaku (perpindahan tab,
keluar dari \emph{fullscreen}, deteksi \emph{restart}, dan berbagai deteksi berbasis perilaku lainnya)

Kata kunci: sistem ujian daring, Computer-Based Testing (CBT), Item Response Theory (IRT), arsitektur SaaS multi-tenant, deteksi kecurangan, sistem konkurensi tinggi
\endabstractind

% % \abstracteng
% % Abstract in English (150-250 words)
%
% Online examination platforms (Computer-Based Testing/CBT) face two critical challenges: maintaining exam integrity against cheating in unsupervised environments, and achieving fair assessment beyond what classical scoring methods can provide. Existing platforms address psychometric scoring or exam integrity separately—none integrate both within a single Software as a Service (SAAS) architecture.
%
% This research aims to design and implement a multi-tenant SAAS architecture for the "Anycademy" examination platform that integrates multi-platform anti-cheat detection with Item Response Theory (IRT)-based assessment. The anti-cheat system employs behavioral detection (tab switching, fullscreen exit, Android screen pinning) with graduated sanctions—temporary block on first and second violations, auto-submit on the third. Scoring uses the 3-Parameter Logistic (3PL) IRT model executed on a Python FastAPI microservice, connected to the Golang modular monolith backend via Redis/Asynq message queue.
%
% Validation was conducted through three types of testing: load and stress testing using K6 to measure throughput and latency under burst traffic scenarios, anti-cheat violation simulations on Android, desktop web, and mobile web platforms to measure detection rate and false positive rate, and interface usability testing using the System Usability Scale (SUS).
%
% Results demonstrate that the proposed architecture handles concurrent exam traffic surges while maintaining exam integrity through multi-platform cheating detection. The combination of asynchronous message queuing and isolated IRT microservices prevents race conditions and thread blocking on the main backend.
%
% Keywords: SaaS architecture, online examination, anti-cheat, Item Response Theory, multi-platform, message queue.
% \endabstracteng

%------------------------------------------------------------
% Chapters
%------------------------------------------------------------
\chapter{PENDAHULUAN}

\section{Latar Belakang}

% TODO: Tulis latar belakang penelitian Anda di sini.
% Jelaskan konteks, urgensi, dan relevansi topik yang dipilih.
% Rujuk pada penelitian terdahulu dan data sekunder jika ada.

Indonesia memiliki lebih dari 316 ribu sekolah yang terkoneksi internet\cite{indonesiabaikRibuSekolah}.
Dari seluruh sekolah tersebut, seluruhnya diwajibkan untuk mengikuti Tes Kemampuan Akademik (TKA) yang berskala nasional 
dan dilakukan menggunakan sarana komputer atau \emph{Computer-Based Testing} (CBT). Hal ini menimbulkan kebutuhan akan 
kegiatan persiapan ujian TKA bagi tiap sekolah. Kegiatan persiapan ujian TKA harus mensimulasikan kondisi 
ujian TKA yang sesungguhnya, hal ini termasuk akan penggunaan software ujian daring yang mirip.

Platform persiapan ujian daring (Computer-Based Testing/CBT) saat ini menghadapi sejumlah tantangan utama untuk digunakan 
sebagai software persiapan atau simulasi bagi institusi. Dua tantangan utama tersebut yaitu ketersediaan algoritma 
penilaian IRT pada software dan kapasitas software untuk menampung dan mengelola data \emph{concurrent} submision 
dalam jumlah besar atau lebih dari 1000 request per detik. Jumlah data sangat penting bagi institusi karena keakuratan 
parameter IRT sangat bergantung pada jumlah data yang tersedia. Hal ini disebabkan oleh adanya potensi penyimpangan 
statistik dari jumlah data yang tersedia akan mempengaruhi keakuratan parameter IRT sehingga menyebabkan
penurunan akurasi skor.

Di sisi lain, tidak semua institusi pendidikan, sekolah, maupun bimbingan belajar memiliki kapabilitas finansial dan
keahlian rekayasa piranti lunak untuk memelihara infrastruktur mandiri demi satu sistem ujian yang pelaksanaannya
bersifat temporal. Kebutuhan akan solusi \textit{hosted} ini mendasari perlunya sebuah perangkat lunak berbiaya rendah,
mudah diakses, dan dapat diandalkan untuk mengelola ujian daring berskala besar.

Penelitian ini menjadikan platform ujian "Anycademy" sebagai studi kasus teknis \emph{software as a service} multi-tenancy.
Arsitektur \textit{backend} akan dirancang menjadi berbasis microservices dengan modul-modul terpisah meliputi 
autentikasi, manajemen dan pemrosesan data tryout, skoring IRT, manajemen \emph{user}, queue, cache, dan database batching.
Fokus penelitian ini adalah pada modul autentikasi, penilaian berbasis IRT, dan pengelolaan data ujian secara berskala besar
dengan simulasi diatas 1000 peserta secara simultan. 

\section{Rumusan Masalah}

% TODO: Tulis rumusan masalah dalam bentuk pertanyaan penelitian.
% Rumusan masalah harus spesifik, terukur, dan dapat dijawab.

Berdasarkan latar belakang di atas, rumusan masalah dalam penelitian ini adalah:

\begin{enumerate}
	\item Bagaimana merancang arsitektur backend platform ujian daring berbasis SaaS yang memisahkan beban proses bisnis umum secara \textit{modular monolith} dengan pemrosesan matematik IRT ke agen layanan mikro (\textit{microservice}) demi menjaga kestabilan platform?
  \item Bagaimana mengembangkan alur pipa (\textit{pipeline}) pendaftaran nilai berantre dengan bantuan \textit{message broker} guna mengelola penyerahan ujian secara konkuren, guna mencegah \textit{race condition} maupun \textit{downtime}?
	\item Bagaimana performa arsitektur yang dikembangkan saat diukur dan divalidasi terhadap skenario \textit{burst traffic} masif dan simulasi pelanggaran anti-kecurangan berdasarkan metrik \textit{throughput}, latensi, dan utilisasi komputasi?
\end{enumerate}

\section{Batasan Masalah}

% TODO: Tulis batasan masalah untuk memfokuskan penelitian.
% Batasan dapat mencakup ruang lingkup, teknologi, atau aspek yang tidak diteliti.

Agar penelitian ini tetap fokus dan terukur di lingkup Ilmu Komputer, batasan masalah yang ditetapkan adalah:

\begin{enumerate}
	\item \textbf{Fokus Arsitektur Backend:} Penelitian difokuskan pada perancangan arsitektur backend platform Anycademy, yang spesifik mencakup sistem \textit{modular monolith} dengan Go, orkestrasi perpesanan dengan Redis, serta integrasi layanan mikro komputasi berat IRT dengan Python FastAPI.
	\item \textbf{Lingkup Eksklusi Psikometrik:} Analisis terhadap sistem IRT dalam sistem ini murni diukur dari performa teknologi komputasionalnya, seperti Big-O ekseskusi, latensi dan efisiensi memori, tanpa mengukur validitas dan ketepatan perhitungan psikometrik secara keilmuan pendidikan.
	\item \textbf{Implementasi Multi-Tenancy:} Metode \textit{multi-tenancy} dibatasi pada implementasi isolasi data secara logis (\textit{Logical Separation/Row-level atau Schema-level}) pada \textit{database} (PostgreSQL), dan tidak mencakup penerapan pemisahan secara fisik (contohnya pembagian klaster basis data bagi masing-masing institusi).
	\item \textbf{Skala Pengujian Beban:} Skenario pengujian skalabilitas (seperti \textit{Load Testing} dan \textit{Stress Testing}) akan mengutilisasi konfigurasi spesifik untuk skenario pengumpulan ujian (*exam submission*) simulasi serentak.
\end{enumerate}

\section{Tujuan Penelitian}
Tujuan dari penelitian ini adalah:
\begin{enumerate}
	\item Mendesain pola arsitektur backend sistem SaaS yang memisahkan antara beban proses bisnis umum secara \textit{modular monolith} dengan pemrosesan matematik IRT ke agen layanan mikro (*microservice*) demi menjaga kestabilan platform.
	\item Mengembangkan alur pipa (*pipeline*) pendaftaran nilai berantre dengan bantuan \textit{message broker} guna mengelola penyerahan ujian secara konkuren, guna mencegah \textit{race condition} maupun \textit{downtime}.
	\item Mengetahui batas puncak performa server (*hardware limit*) beserta skalabilitas piranti lunaknya lewat analisis statistik kuantitatif terhadap \textit{throughput} (T/s), waktu respons (\textit{response latency}), dan penggunaan CPU/RAM selama simulasi ujian intensif divalidasi.
\end{enumerate}

\section{Manfaat Penelitian}

Hasil dari penelitian ini diharapkan dapat memberikan kontribusi sebagai berikut:

\begin{enumerate}
	\item \textbf{Manfaat Teoritis:} Penelitian ini berkontribusi pada literatur rekayasa perangkat lunak dengan menyajikan
  studi kasus konkret mengenai pola arsitektur hibrida (\textit{polyglot architecture}) untuk domain ed-tech yang 
  menuntut komputasi intensif sekaligus ketahanan terhadap lonjakan trafik konkuren. Secara khusus, penelitian ini 
  memperkaya pemahaman tentang bagaimana pemisahan tanggung jawab antara proses bisnis umum (\textit{modular monolith}) 
  dan komputasi numerik berat (\textit{microservice}) dapat dirumuskan sebagai pola desain yang dapat direplikasi pada 
  sistem serupa.

	\item \textbf{Manfaat Praktis (Industri Perangkat Lunak \& Pengembang EdTech):} Penelitian ini diharapkan dapat menjadi 
  rujukan teknis dan cetak biru (\textit{blueprint}) arsitektur bagi pengembang maupun tim rekayasa yang merancang sistem 
  ujian daring berskala besar, khususnya dalam mengombinasikan keunggulan konkurensi bahasa Go untuk lapisan layanan 
  inti dengan kekuatan ekosistem numerik/ilmiah Python untuk beban komputasi psikometrik. Dengan demikian, hasil penelitian 
  dapat mempercepat proses pengambilan keputusan arsitektural pada proyek-proyek ed-tech serupa, tanpa perlu melalui proses 
  trial-and-error dari nol.

\end{enumerate}


\section{Sistematika Penulisan}

% TODO: Sesuaikan sistematika penulisan dengan struktur penelitian Anda.

Sistematika penulisan proposal tugas akhir ini adalah sebagai berikut:

\begin{itemize}
	\item \textbf{Bab I:} Pendahuluan --- berisi pendefinisian latar belakang problem rekayasa sistem, rumusan arsitektural, ruang lingkup implementasi teknis, luaran fungsionalitas, serta manfaat komputasional program.
	\item \textbf{Bab II:} Tinjauan Pustaka --- merangkum karya-karya terkait \textit{software scaling}, perbandingan arsitektur (\textit{Monolith vs Microservices}), manajemen antrean asinkronus, serta teknologi landasan CBT.
	\item \textbf{Bab III:} Metode Penelitian --- menjabarkan alur metodologi rekayasa perancangan arsitektur, diagaram topologi \textit{SaaS}, pemodelan integrasi modul, skema basis data, dan rancangan pengujian \textit{stress \& load test}.
	\item \textbf{Bab IV:} Jadwal Penelitian --- memuat Gantt Chart linimasa fase perancangan hingga implementasi prototipe.
\end{itemize}

%\chapter{TINJAUAN PUSTAKA}

% Literature review covering theoretical foundations and practical applications
% of scalable software architecture, distributed systems for CBT, and ISO-standardized UI/UX.

\section{Komputasi Teori Respons Butir (IRT)}

Teori Respons Butir (Item Response Theory, IRT) adalah kerangka psikometrik yang membangun asumsi bahwa kinerja responden
pada suatu butir soal dipengaruhi oleh sifat-sifat laten yang dimiliki responden (kemampuan atau ability, $\theta$) 
dan parameter-parameter butir yang spesifik \citep{Hambleton1985, Baker2004}. Model logistik parametrik seperti 
1-PL, 2-PL, dan 3-PL telah menjadi standar dalam pengujian global maupun nasional seperti UTBK.

Secara matematis, probabilitas responden $i$ menjawab butir soal $j$ dengan benar dalam model 3-PL didefinisikan sebagai:

\begin{equation}
	P_{ij}(\theta_i) = c_j + (1 - c_j) \frac{e^{a_j(\theta_i - b_j)}}{1 + e^{a_j(\theta_i - b_j)}}
\end{equation}

dimana $a_j$ adalah parameter diskriminasi, $b_j$ adalah kesulitan butir, dan $c_j$ adalah paramater tebakan 
\citep{Birnbaum1968}. Dari perspektif Ilmu Komputer dan Rekayasa Perangkat Lunak, kalibrasi parameter IRT secara dinamis 
untuk ribuan siswa secara konkuren adalah proses komputasional yang intensif (\textit{CPU-bound process}). 
Estimasi parameter seperti \textit{Maximum Likelihood Estimation} (MLE) atau Pendekatan Bayesian melibatkan komputasi 
matriks berulang dan iterasi numerik yang membebani unit pemrosesan sentral. Hal ini membatasi eksekusi algoritma ini 
pada bahasa pemrogaman yang tidak dioptimalkan untuk komputasi saintifik. Oleh karenanya, layanan pengolahan IRT umumnya 
diisolasi ke dalam kerangka layanan mikro (\textit{microservices}) berbasis Python (memanfaatkan optimasi 
librari operasi pelarik tingkat level C seperti NumPy maupun SciPy) daripada dijalankan pada utas peladen 
transaksional (*routine*) layaknya fungsi CRUD berbasis Golang.

\section{Sistem Ujian Berbasis Komputer \& Lonjakan Lalu Lintas (Burst Traffic)}

Computer-Based Testing (CBT) menggeser ujian tradisional berbasis kertas menuju media digital. Salah satu tantangan 
rekayasa infrastruktur yang paling mendasar dalam operasional CBT berskala besar adalah fenomena \textit{thundering herd} 
atau lonjakan lalu lintas yang terpusat dan berdurasi pendek (\textit{burst traffic}). Dalam skenario ujian dunia nyata, 
ratusan atau bahkan ribuan siswa pada sekolah yang sama akan mengakses sistem dan melakukan pengumpulan jawaban 
(\textit{exam submission}) secara serentak pada hitungan detik yang sama saat bel hitung mundur waktu ujian berakhir. 
Lonjakan konkurensi ini memicu pelipatgandaan HTTP \textit{request} secara seketika, yang pada gilirannya menyebabkan 
interupsi peladen (\textit{downtime}), hambatan pita jaringan (\textit{network bottleneck}), dan penipisan jumlah koneksi 
pangkalan data (\textit{database connection pool exhaustion}) pada arsitektur \textit{monolithic} tradisional. Penanganan 
\textit{burst traffic} membutuhkan desain konkurensi tingkat lanjut menggunakan pendekatan pemrograman reaktif atau 
arsitektur antrean asinkron agar terhindar dari pemblokiran utas eksekusi (\textit{thread blocking}) dan kondisi balapan 
(\textit{race condition}).

\section{Integritas Ujian dan Metode Anti-Kecurangan dalam CBT}

Salah satu kelemahan fundamental Computer-Based Testing adalah kerentanan terhadap kecurangan akibat ketiadaan pengawas 
fisik \citep{Wibowo2018,Sorensen2018}. Metode anti-kecurangan dalam CBT dapat diklasifikasikan menjadi empat kategori:

\begin{enumerate}
	\item \textbf{Deteksi Berbasis Perilaku (\textit{Behavior-Based}):} Mendeteksi perubahan status perangkat seperti perpindahan tab peramban, keluar dari mode layar penuh, aktivasi \textit{screen saver}, atau restart perangkat. Keunggulannya adalah non-invasif dan kompatibel lintas platform \citep{Dawson2017}.

	\item \textbf{Pengawasan Berbasis Biometrik (\textit{Biometric Proctoring}):} Menggunakan webcam untuk mendeteksi gerakan mata (\textit{gaze tracking}), ekspresi wajah, dan keberadaan orang lain di ruangan. Metode ini akurat namun invasif dan membutuhkan bandwidth tinggi \citep{Chuang2020}.

	\item \textbf{Deteksi Berbasis Analitik (\textit{Analytic-Based}):} Menganalisis pola jawaban untuk mendeteksi kemiripan antar peserta (\textit{answer copying}), waktu pengerjaan yang tidak wajar, atau pola jawaban yang mengindikasikan kerja sama \citep{Wollack2003}.

	\item \textbf{Pengendalian Lingkungan (\textit{Environment Control}):} Membatasi akses sistem operasi melalui \textit{lockdown browser} (seperti Safe Exam Browser, Respondus) atau \textit{Android screen pinning} (LockTask mode) yang mengunci perangkat ke satu aplikasi \citep{Franchi2019}.
\end{enumerate}

Platform ujian komersial yang ada menerapkan metode-metode di atas secara terpisah dan tidak terintegrasi dengan sistem penilaian psikometrik. Tabel \ref{tab:platform-comparison} menunjukkan perbandingan fitur antar platform.

\begin{table}[H]
	\centering
	\small
	\caption{Perbandingan Platform Ujian Online}
	\label{tab:platform-comparison}
	\begin{tabular}{|l|c|c|c|c|c|}
		\hline
		\textbf{Fitur}            & \textbf{Anycademy} & \textbf{Moodle} & \textbf{ProctorU} & \textbf{Exam.net} & \textbf{G-Forms} \\
		\hline
		Penilaian IRT             & $\bullet$          & --              & --                & --                & --               \\
		\hline
		Anti-cheat multi-platform & $\bullet$          & $\sim$          & $\bullet$         & $\bullet$         & --               \\
		\hline
		Eskalasi berjenjang       & $\bullet$          & --              & --                & --                & --               \\
		\hline
		SAAS multi-tenant         & $\bullet$          & $\bullet$       & $\bullet$         & $\bullet$         & $\bullet$        \\
		\hline
		Async message queue       & $\bullet$          & --              & --                & --                & --               \\
		\hline
		AI PDF parsing soal       & $\bullet$          & --              & --                & --                & --               \\
		\hline
		Bisnis Indonesia          & $\bullet$          & $\bullet$       & --                & --                & $\bullet$        \\
		\hline
	\end{tabular}
\end{table}

Keterangan: $\bullet$ = Tersedia; $\sim$ = Terbatas (plug-in); -- = Tidak tersedia.

Dalam konteks penelitian ini, Anycademy mengadopsi pendekatan \textbf{deteksi berbasis perilaku} (kategori 1) yang dikombinasikan dengan \textbf{pengendalian lingkungan} (kategori 4) — yaitu Android \textit{screen pinning} via Capacitor, \textit{fullscreen enforcement} pada browser, serta deteksi \textit{tab-switch} dan \textit{device restart}. Mekanisme eskalasi hukuman berjenjang (\textit{graduated sanction}) diterapkan: pelanggaran pertama mengakibatkan pemblokiran sementara (10 menit), pelanggaran kedua memperpanjang blokir (20 menit), dan pelanggaran ketiga memicu pengumpulan otomatis (\textit{auto-submit}) ujian.

\section{Arsitektur Perangkat Lunak Terskala Tinggi}

Untuk mengatasi keterbatasan stabilitas perangkat lunak, paradigma pola desain arsitektur berskala tinggi mutlak diaplikasikan:

\subsection{Microservices versus Monolithic}

Arsitektur sinkronus tradisional \textit{Monolithic} menyatukan seluruh antarmuka, logika bisnis, dan pengolahan basis data dalam satu entitas layanan \citep{Newman2015}. Kelebihannya terletak pada kemudahan implementasi pemantauan jejak pelacakan galat, namun skalabilitasnya amatlah terbatas dan rentan ambruk saat sistem harus mengkompensasi beban asimetris antar fiturnya.

Sebagai antitesis, arsitektur \textit{Microservices} memecah fungsionalitas logika aplikasi menjadi layanan-layanan portabel dan tersendiri \citep{Newman2015}. Metode agnostik perantara ini memastikan bahwa jika fitur layanan penskoran matematis IRT mendadak menerima utilisasi performa ekstrim hingga batas limitnya, maka rekayasawan cukup memberhentikan atau menggandakan klaster peladen layanan IRT itu saja (\textit{scale-up}/\textit{scale-out}), alih-alih memberatkan modul transaksi ujian lain. Dalam konteks Anycademy, metode perpaduan (hybrid) diimplementasikan; sebuah \textit{Modular Monolith} (berbasis sistem Go)\ dipertahankan lantaran kecepatannya mendistribusikan modul manajemen *state* dan ruter, disinkronisasikan bersama \textit{Rest API microservices} cerdas kecerdasan buatan dan perhitungan kalkulasi IRT mandiri menggunakan perantara pustaka Python FastAPI.

\section{Sistem Message Broker dan Antrean Asinkron (Async Queue)}

Strategi optimasi komputasi paling krusial demi mengeliminasi penumpukan beban kueri (*overload*) selama lonjakan ujian massal adalah pelolosan arsitektur *asynchronous* melalui peran mediator pialang pesan (*message broker*) digabung skema alur distribusi antrean (*distributed task queues*).

Penggunaan pangkalan struktur penyimpanan berbasis ingatan utama kilat (*in-memory*) seperti Redis bersama dengan perpustakaan sinkronisasi latar belakang seperti Asynq, memfasilitasi arsitektur pola "Penerbit-Pelanggan" (\textit{Publisher-Subscriber} / Pub-Sub). Ketimbang membiasakan proses penyimpanan kueri jawaban dari belasan ribu terminal ujian membebani baris tulis transaksional PostgreSQL secara sinkron dan mendempet jeda kompilasinya, *web backend* diarahkan sekedar meresahkan persetujuan HTTP 202 ("Jawaban Diterima dalam Antrean"), mendelegasikan payload muatan skoring beruntun, selagi \textit{background worker task} diam-diam me-resolve kalkulasi probabilitas matriks algoritmik IRT dan persistensi penulisan basis data belakangan guna murni mengakomodir kenyamanan layar interaksi peserta tes yang steril (\textit{seamless execution}).

\section{Multi-Tenant SaaS dan Isolasi Data}

Dari sisi komersialisasinya, sistem Anycademy diusung terarsitektur melalui wujud integrasi \textit{Software/System as a Service} (SaaS) bermodel entitas jamak (\textit{multi-tenant}). Multi-tenancy melingkupi metode khusus rekayasa penyebaran di mana sebentuk (\textit{single-instance}) klon aplikasi kode sumber tunggal serta hierarki penyetelan model peladen pangkalan datanya dikonfigurasi simultan demi menyokong berganda penyewa terpisah (klien sekolah, lembaga bimbingan mengajar, institusi, dst.).

Risiko kelemahan dari platform ed-tech tersentralisasi sedemikian adalah terobosan penyekatan privasi dan intervensi batas otoritas rekam jejak (\textit{Security Privileges}). Isolasi logika perpaduan lintas kolom (\textit{Logical Isolation / Shared Database but Isolated Schema}) mereduksi kekhawatiran itu pada beban infrastruktur termurah secara holistik (\textit{Cost Effective}). Validasi restriktif sistem manajemen data RDBMS pada peladen PostgreSQL dipertegas pada lapis regulasi \textit{Row-Level Security} (RLS). Setiap akses terotentikasi dimitigasi secara silang dengan referensi identitas Organisasi/Tenant terkekang yang disusupkan lewat otorisasi berlapis keping \textit{JSON Web Token} (JWT) \citep{Khan2013}.

\section{Lingkaran Interaksi Manusia-Komputer dengan User-Centered Design (ISO 9241-210)}

Dalam industri ed-tech lintas skala SaaS, orkestrasi *backend server* solid tidak akan optimal perancangannya, bilamana dihadapkan kepada kegagalan usabilitas pengguna (*human error*) dikarenakan rancangan antarmuka yang mengacaukan limit memori kognitif klien awam (\textit{cognitive overload}). Pendekatan empiris disiplin teknik Interaksi Manusia-Komputer (HCI) difasilitasi ke dalam kaidah rancang bangun UCD berpedoman standarisasi fabrikasi internasional, yaitu ISO 9241-210: *Human-centred design for interactive systems* \citep{ISO92412010}.

\subsection{Rekayasa Antarmuka UI/UX Berpusat Pengguna}

Model tahapan pengembangan ISO merekomendasikan literasi desain iteratif dan siklus usabilitas *prototyping* kualitatif berbasis kebutuhan langsung target *user*. Di area sistem CBT, implementasi lini layar sentuh UI dan alur pengalaman UX bercabang menjadi dua rute fungsionalitas tumpu:

\begin{enumerate}
	\item \textbf{CBT Interaktif bagi Siswa (Laman Ujian Asesmen):} Lingkungan antarmuka yang presisi direkayasa demi meminimalisasi redudansi letak opsi fungsi pemicu, bersih distorsi tata ruang, hingga keandalan koneksi navigasi *autosave* di atas gejolak jaringan lemah, untuk meredakan limit beban stres saat siswa menghadapi tensi tempo ujian penentu tinggi (*High-Stakes Session*).
	\item \textbf{Dasbor Dasbor Analitik Pendidik (Teacher Portal):} Transkripsi teknis grafis kalkulasi kompleks eksponen dan diskriminasi IRT didekonsolidasikan, mensimplifikasi indikator logistik logaritma mentah menuju panel \textit{insight} wawasan representasi parameter siswa yang lebih intuitif namun esensial merangkum efektivitas *actionable learning decision* pengajaran.
\end{enumerate}

\subsection{Kuantifikasi Validasi System Usability Scale (SUS)}

Menyoal validasi teknikal pasca konstruksi prototipe, verifikasi tingkat kompatibilitas keberhasilan sasaran desain layar dari ISO 9241-210 (\textit{Usability Testing}) ditranslasikan dalam kuantifikasi metrik industri bernilai matematis, yang dikarakteristik melalui tes penyelesaian \textit{System Usability Scale} (SUS) \citep{Brooke1996}. Instrumen pengukuran koefisien SUS terdiri atas sepuluh rangkaian pernyataan umpan survei kuisioner \textit{likelihood} skor 1 hingga 5, mensitesiskan efektivitas kemudahan orientasi antarmuka sasaran; yang lazim dimutlakkan sukses bernilai positif bilamana rekam skor rata rata subjek penilai minimal konsisten meraih klasifikasi metrik akhir $\geq 68$.

\section{Pengujian Validitas Skalabilitas Kinerja Perangkat Lunak}

Kesiapan unjuk gigi performa implementasi infrastruktur \textit{backend} akan divisualisasikan melalui tes stres simulatur komprehensif mengacu kepada parameter audit optimisasi perangkat lunak awam:
\begin{itemize}
	\item \textbf{Throughput:} Frekuensi daya penyelesaian akumulasi paket \textit{I/O requests/transactions} bersih yang sukses diproses peladen hitungan persatuan selang waktu detik — identik didefinisikan sebagai rasio \textit{Requests Per Second} (RPS).
	\item \textbf{Latensi (Latency/Response Response Time):} Metrik konversi sekuen temporal delta waktu rekam jejak kapan komando antarmuka API terminal pertama kali direquest (\textit{send}) hingga *payload content* usai dikirim kembali utuh (\textit{acknowledge response}).
	\item \textbf{Tingkat Eksepsi Trafik (Error Rate):} Angka rasio probabilitas kegagalan eksepsi jaringan serta I/O \textit{fatal bug}, yang memetakan tumpahan interupsi \textit{HTTP 5xx Server Error} hingga frekuensi sambungan pemutusan waktu tenggang (\textit{connection timeouts}) akibat peladen tidak kuat menyanggupi lonjakan paksaan beban dalam interval padat serentak.
	\item \textbf{Utilisasi Komputasi (Resource Metrics):} Batas parameter pemrosesan penggunaan inti sirkuit pemroses data prosesor terunggah (\textit{CPU Allocation}) maupun daya penyimpanan memori \textit{Memory/RAM}.
\end{itemize}
Metodologi tinjauan penjelajahan audit peranti di atas divalidasi presisi memakai perangkat bantu replikator lonjakan rekayasa kueri maya bernama \textit{Load Testing} (simulasi perilaku padat pengguna paralel skala riil) dilanjutkan pengerasan kompresi beruntun lewat \textit{Stress Testing} (penentuan *breaking/crash points* saat rentetan lalu lintas memprovokasi muatan paksaan abnormalitas trafik melampaui batasan kapasitas kognitif mesin arsitektur yang disediakan).

\section{Kesenjangan Penelitian (Research Gap) dan Pemosisian}

Berdasarkan tinjauan pustaka di atas, tiga kesenjangan penelitian (\textit{research gap}) teridentifikasi:
\begin{enumerate}
	\item Platform ujian online yang ada menangani penilaian psikometrik (IRT) \textbf{atau} integritas ujian (anti-cheat), tetapi belum ada yang mengintegrasikan keduanya dalam satu arsitektur SAAS.
	\item Mekanisme anti-kecurangan yang ada umumnya bersifat biner (\textit{block}/\textit{allow}) tanpa eskalasi hukuman berjenjang yang memberikan kesempatan koreksi kepada peserta.
	\item Literatur yang mendokumentasikan pengujian beban arsitektur SAAS ujian yang menggabungkan \textit{message queue}, \textit{microservices} IRT, dan deteksi anti-kecurangan secara simultan masih terbatas.
\end{enumerate}

Penelitian ini memposisikan Anycademy sebagai solusi yang mengisi ketiga kesenjangan tersebut: sebuah platform SAAS CBT yang mengintegrasikan deteksi kecurangan multi-platform dengan eskalasi hukuman berjenjang dan penilaian IRT dalam arsitektur \textit{modular monolith} + \textit{microservices}. Validasi dilakukan melalui pengujian beban sintetik (\textit{load testing} dengan K6) dan pengujian efektivitas deteksi anti-kecurangan pada tiga platform target (Android, desktop browser, dan web).

%\include{chapters/bab3-metodologi}
%\chapter{JADWAL PENELITIAN}

Penelitian ini direncanakan berlangsung selama \textbf{4 bulan} dengan fokus pada perancangan, implementasi, pengujian, dan dokumentasi arsitektur perangkat lunak Anycademy. Jadwal disusun berdasarkan alur rekayasa perangkat lunak yang dijelaskan pada Bab III.

\section{Ringkasan Fase}

\begin{enumerate}
	\item \textbf{Fase 1 --- Perancangan Arsitektur (Bulan 1):} Analisis kebutuhan sistem, perancangan topologi arsitektur (\textit{modular monolith} + \textit{microservices}), perancangan skema basis data (PostgreSQL, MongoDB, Redis), dan perancangan API endpoints.

	\item \textbf{Fase 2 --- Implementasi (Bulan 1--3):} Pengembangan backend (Go/Gin), \textit{microservices} (Python FastAPI: IRT + AI), frontend (React/Vite), integrasi Redis/Asynq \textit{message queue}, implementasi \textit{multi-platform anti-cheat}, dan \textit{deployment} via Docker Compose.

	\item \textbf{Fase 3 --- Pengujian (Bulan 3--4):} \textit{Load} dan \textit{stress testing} dengan K6, pengujian efektivitas anti-kecurangan pada tiga platform, pengujian usabilitas antarmuka dengan \textit{System Usability Scale} (SUS), dan evaluasi performa \textit{pipeline} IRT.

	\item \textbf{Fase 4 --- Dokumentasi \& Pelaporan (Bulan 3--4):} Penulisan dokumentasi teknis, analisis hasil pengujian, penyusunan laporan tugas akhir, dan finalisasi naskah.
\end{enumerate}

\section{Gantt Chart}

\begin{table}[H]
	\centering
	\caption{Jadwal Penelitian Anycademy --- 4 Bulan}
	\begin{tabular}{|l|c|c|c|c|}
		\hline
		\multirow{2}{*}{\textbf{Kegiatan}}     & \multicolumn{4}{c|}{\textbf{Bulan}}                                        \\
		\cline{2-5}
		                                       & \textbf{1}                          & \textbf{2} & \textbf{3} & \textbf{4} \\
		\hline
		\multicolumn{5}{|c|}{\textbf{Fase 1: Perancangan}}                                                                  \\
		\hline
		Studi literatur \& analisis kebutuhan  & $\bullet$                           &            &            &            \\
		\hline
		Perancangan arsitektur sistem          & $\bullet$                           &            &            &            \\
		\hline
		Perancangan skema basis data           & $\bullet$                           &            &            &            \\
		\hline
		Perancangan API endpoints              & $\bullet$                           & $\bullet$  &            &            \\
		\hline
		\hline
		\multicolumn{5}{|c|}{\textbf{Fase 2: Implementasi}}                                                                 \\
		\hline
		Pengembangan backend (Go/Gin)          & $\bullet$                           & $\bullet$  & $\bullet$  &            \\
		\hline
		Pengembangan IRT Service (Python)      & $\bullet$                           & $\bullet$  &            &            \\
		\hline
		Pengembangan AI Service (Python)       & $\bullet$                           & $\bullet$  &            &            \\
		\hline
		Pengembangan frontend (React/Vite)     & $\bullet$                           & $\bullet$  & $\bullet$  &            \\
		\hline
		Integrasi Redis/Asynq queue            &                                     & $\bullet$  & $\bullet$  &            \\
		\hline
		Implementasi multi-platform anti-cheat &                                     & $\bullet$  & $\bullet$  &            \\
		\hline
		Deployment Docker Compose              &                                     & $\bullet$  & $\bullet$  &            \\
		\hline
		\hline
		\multicolumn{5}{|c|}{\textbf{Fase 3: Pengujian}}                                                                    \\
		\hline
		Load \& stress testing (K6)            &                                     &            & $\bullet$  & $\bullet$  \\
		\hline
		Pengujian efektivitas anti-cheat       &                                     &            & $\bullet$  & $\bullet$  \\
		\hline
		Pengujian usabilitas (SUS)             &                                     &            &            & $\bullet$  \\
		\hline
		Evaluasi performa IRT pipeline         &                                     &            & $\bullet$  &            \\
		\hline
		\hline
		\multicolumn{5}{|c|}{\textbf{Fase 4: Dokumentasi \& Pelaporan}}                                                     \\
		\hline
		Penulisan dokumentasi teknis           &                                     & $\bullet$  & $\bullet$  & $\bullet$  \\
		\hline
		Analisis hasil pengujian               &                                     &            & $\bullet$  & $\bullet$  \\
		\hline
		Penyusunan laporan TA                  &                                     &            & $\bullet$  & $\bullet$  \\
		\hline
		Finalisasi naskah                      &                                     &            &            & $\bullet$  \\
		\hline
	\end{tabular}
\end{table}

\noindent \textbf{Keterangan:} $\bullet$ = Kegiatan dilaksanakan pada bulan tersebut.

\section{Milestone Kritis}

\begin{itemize}
	\item \textbf{Akhir Bulan 1:} Desain arsitektur dan skema basis data selesai. Backend \textit{core} (auth, exam CRUD) berfungsi di lingkungan lokal.

	\item \textbf{Akhir Bulan 2:} Backend dan frontend terintegrasi penuh. IRT Service dan AI Service berfungsi. Sistem anti-cheat (\textit{KIOSK mode}, \textit{tamper detection}) terimplementasi pada web dan Android.

	\item \textbf{Akhir Bulan 3:} Seluruh fitur terimplementasi. Pengujian beban dengan K6 dan pengujian efektivitas anti-cheat selesai. Hasil pengujian terdokumentasi.

	\item \textbf{Akhir Bulan 4:} Laporan tugas akhir selesai, termasuk analisis hasil pengujian, kesimpulan, dan saran pengembangan. Naskah siap diserahkan.
\end{itemize}


%------------------------------------------------------------
% Bibliography
%------------------------------------------------------------
\bibliography{references}

\end{document}
